\documentclass[11pt]{article}
\usepackage[margin=0.9in]{geometry}
\usepackage{amsmath,amssymb,amsthm}
\newtheorem{lemma}{Lemma}
\newcommand{\R}{\mathbb R}
\title{Necessary face and contact conditions for a ramified selector}
\author{}\date{}
\begin{document}\maketitle\vspace{-2em}
Use the canonical source order
$c=(d,\lambda,n,m,C,A,D,\alpha,\beta)\ge0$ and the matrices $P,N,E,K$
of the canonical-completion note. Put $v=16s^2$ and let
$J(s)\in(\R[s]_{\le4})^{2\times2}$. Write $R=K-E-J$ and
\[
 G=\begin{pmatrix}P+K&J\\J^T&N+K\end{pmatrix},\qquad
 H=\begin{pmatrix}P+E&R\\R^T&N+E\end{pmatrix}.
\]
The nine source coordinates remain independent. The conditions below
are necessary for positivity on $s\ge0$ and do not suffice for a completion.

\begin{lemma}[Endpoint support and contact interpolation]\label{lem:necessary}
Let $D_i^M(s)$ be diagonal entry $i$ of $M=G,H$. With $\delta=d/2$,
$D_i^M(s)=4(a_i^M+256b_i^Ms^4)$, where
\[
\begin{array}{c|cc|cc}
 i&a_i^G&b_i^G&a_i^H&b_i^H\\\hline
 1&\lambda&m&A&m+A+\alpha+\delta\\
 2&m+\beta+n+\delta&\beta+\lambda+n+\delta
   &m+\beta+\alpha+\delta&\beta\\
 3&C+D+\lambda+\delta&C&C+D+A+\delta&C+A+\alpha+\delta\\
 4&n&D+\lambda+n+\delta&\alpha&D
\end{array}
\]
For $i,j\in\{1,2\}$, an entry $Z_{ij}$, with $Z=J$ for $G$ and
$Z=R$ for $H$, vanishes identically if either $D_i^M$ or $D_{j+2}^M$
does. Otherwise every exponent $k$ occurring in it satisfies
\begin{equation}
 \frac{\operatorname{ord}_0D_i^M+\operatorname{ord}_0D_{j+2}^M}{2}
 \le k\le
 \frac{\deg D_i^M+\deg D_{j+2}^M}{2}.                \label{eq:support}
\end{equation}
Here the orders and degrees of a nonzero diagonal are $0$ or $4$.
Thus these conditions are explicit affine restrictions on the twenty
coefficients of $J$, on every source face.

At an interior full-state contact $(U_0,s_0;a,b)$, where $U_0>0$,
$s_0>0$, and the canonical form $\Phi_c(U_0,16s_0^2;a,b)=0$, put
$g=(a,U_0b)^T$, $h=(a,b)^T$. Then
\begin{align}
 G(s_0)g&=0,& H(s_0)h&=0,\notag\\
 J(s_0)b&=-(P+K)(s_0)a/U_0,&
 J(s_0)^Ta&=-U_0(N+K)(s_0)b,                       \label{eq:values}\\
 2U_0a^TJ'(s_0)b
 &=-a^T(P+K)'(s_0)a-U_0^2b^T(N+K)'(s_0)b.         \label{eq:derivative}
\end{align}
Primes mean differentiation with respect to $s$. These formulas require
no positive source coefficient and no nonzero state component.
\end{lemma}
\begin{proof}
The displayed diagonals follow by direct substitution; in particular
the signed $d$ terms cancel in the leading coefficients of $H_{22}$
and $H_{44}$. All table entries are nonnegative.
Every $2\times2$ principal minor gives
$|Z_{ij}(s)|^2\le D_i^M(s)D_{j+2}^M(s)$.
The lowest and highest powers as $s\downarrow0$ and $s\to\infty$
prove \eqref{eq:support}, including all literal zero cases.

The exact identity
\[
 \Phi_c(U,16s^2;a,b)=
 \binom a{Ub}^{\!T}G(s)\binom a{Ub}
 +U\binom ab^{\!T}H(s)\binom ab
\]
makes both terms zero at a contact, so each PSD matrix kills its
displayed vector. This gives \eqref{eq:values}.
For a fixed kernel vector, the corresponding scalar quadratic form
has a minimum at the interior point $s_0$; differentiating $g^TG(s)g$
gives \eqref{eq:derivative}.
The analogous derivative from $H$ is consistent: its sum with the
$G$ derivative, weighted by $U_0$, is $\partial_s\Phi_c=0$.
\end{proof}

On C1, $\lambda=n=0$, $d,m,C>0$, the first support restriction gives
\begin{equation}
 J(s)=\begin{pmatrix}
 p_2s^2+p_3s^3+p_4s^4&q_4s^4\\
 r_0+r_1s+r_2s^2+r_3s^3+r_4s^4&\ell_2s^2+\ell_3s^3+\ell_4s^4
 \end{pmatrix}.                                    \label{eq:c1-form}
\end{equation}
Only four coefficients are added to an unramified quadratic selector:
$p_3,\ell_3,r_1,r_3$. The second support restriction also retains
the literal source faces. If $D=0$, necessarily
$J_{12}=0$ and $J_{22}=\ell_2s^2+512ds^4$.
If $\beta=0$, necessarily $\deg J_{21}\le2$ and
$J_{22}=\ell_2s^2+512ds^4$.
If both vanish, $J_{22}=512ds^4$.
Additional zero sources are handled directly by the table; no passage
to a strictly positive source is needed.

\paragraph{The remaining local positivity condition.}
The linear contact equations do not establish even local matrix
positivity. For either $M=G,H$ at an interior contact, let $V$ span
$\ker M(s_0)$ and let $T$ span a complementary subspace. Then
$D=T^TM(s_0)T\succ0$ and local positivity necessarily implies
\begin{equation}
 V^TM'(s_0)V=0,\qquad
 \frac12V^TM''(s_0)V
 -(T^TM'(s_0)V)^TD^{-1}(T^TM'(s_0)V)\succeq0.       \label{eq:schur-jet}
\end{equation}
This follows by expanding the Schur complement to second order.
If the complementary space is empty, the inverse term is zero.
At $s=0$ only the one-sided condition $V^TM'(0)V\succeq0$ is
automatic; if that matrix vanishes, the second condition above applies.
The same boundary rule holds at infinity after reversing the polynomial
as $\tau^4M(1/\tau)$.

The physical kernel compatibility makes the value equations exact:
if $\mathcal P_c h=0$ for the state matrix of $\Phi_c$, then $Gg=0$
implies $Hh=0$ by its matrix splitting identity.
When $a,b\ne0$, the value equations impose three independent linear
conditions. The difference of any two solutions is exactly
$t\,a_\perp b_\perp^T$, where $a_\perp=(-a_2,a_1)^T$ and similarly
for $b$. If $a=0,b\ne0$, the homogeneous freedom is
$w b_\perp^T$; if $b=0,a\ne0$, it is $a_\perp w^T$, with $w\in\R^2$.
The prescribed right sides must of course be compatible.
The first jets impose scalar kernel expectations, not $G'g=0$ or
$H'h=0$. The same twenty polynomial
coefficients must satisfy these conditions at every contact, the two
endpoint support restrictions, and the coupled Schur inequalities
\eqref{eq:schur-jet}. Rank drops use the full kernel $V$, rather than
only the physical contact vector. Even these second-order conditions
are not sufficient: $-(s-s_0)^4$ passes the corresponding scalar
zero, first- and second-derivative tests but is negative nearby.

The unresolved quantifier is therefore still
\[
 \forall c\in C_{\rm can}\quad\exists J_0,\ldots,J_4\in\R^{2\times2}
 \quad\forall s\ge0:\quad G(s),H(s)\succeq0.
\]
Pointwise completion fixes neither the required derivatives nor a common
quartic interpolant. Factorization of an already positive quartic
matrix also does not supply that interpolant. The lemma gives necessary
conditions for a universal proof; it proves neither universal adequacy
nor a counterexample to it.
\end{document}
